%%%%%%%%%%%%%%%%%%%%%%%%%%%%%%%%%%%%%%
\section{Systém DNS (Domain Name System)}\label{dns}
Systém DNS slouží jako podpora internetové komunikace, kdy ukládá informace důležité ke správnému fungování počítačových sítí a služeb. Jednou z nejvíce využívaných služeb je překlad mezi doménovým jménem na IP adresu a naopak, tzn. vyhledávání záznamů typu A, AAAA a PTR \cite{rfc1035,rfc3596}.

Mezi další důležité záznamy patří záznam {\tt MX}, který obsahuje odkaz na adresu poštovního serveru pro danou doménu. Záznamy {\tt MX} jsou nezbytné pro správné fungování systému elektronické pošty nad SMTP. Pro vyhledávání serverů dalších služeb (např. SIP, XMPP a další) slouží záznamy typu typu {\tt SRV}, které lze v~současnosti využít pro služby SIP a XMPP \cite{rfc2782}.

Některé záznamy během vývoje internetu rozšířili svou funkci. Například záznamy typu {\tt TXT} byly původně používány pro uložení textových informací o dané doméně. V současné době se převážně využívají pro zajištění bezpečnosti softwarových služeb či přenosu elektronické pošty, kde záznam {\tt TXT} například obsahuje veřejný klíč pro službu DKIM (DomainKeys Identified Mail) \cite{rfc6376}).

Služba DNS pracuje na principu klient--server, kde klient (zvaný též resolver) posílá dotaz pro vyhledání informací v distribuované globálním systému DNS. Dotaz přijde na server DNS, který se snaží najít odpověď bud v lokální databázi (zónovém souboru) či paměti cache. Pokud odpověď nenajde lokálně, může se dotazovat dalších serverů (v případě rekurzivního dotazu). Proces vyhledání informací v systému DNS se nazývá {\em rezoluce}. 

\subsection{Konfigurace klienta DNS}
Co je klient DNS? Možná zvláštní otázka, nicméně definice klienta DNS není v počítačových sítích zcela jednoznačná. Za klientské programy lze považovat aplikace typu {\tt nslookup}, {\tt host} či {\tt dig}, což jsou programy, které se dotazují systému DNS. Všechny tyto programy pro svou činnost využívají soubor systémových rutin zvaný {\em resolver}, který je součástí operačního systému a implementuje zasílání dotazů na server DNS a zpracování odpovědí. Z uživatelského pohledu lze za klienty DNS považovat výše uvedené aplikace typu {\tt nslookup}, z implementačního pohledu je klientem DNS {\em resolver}. Protože však aplikační programy typu {\tt nslookup} umožňují uživateli přistupovat k systému DNS, tzn. komunikují se servery DNS, z formální pohledu je v pořádku, když je také budeme nazývat klienty DNS.

Pro konfiguraci klienta DNS se využívají následující konfigurační soubory: {\tt /etc/hosts}, {\tt /etc/re\-solv.conf} a {\tt /etc/hosts.conf}. 
\begin{itemize}
  \item {\tt /etc/hosts} -- Tento soubor obsahuje statické mapování IP adres a doménových jmen. Vytváří se manuálně a využívá se například jako záloha v situaci, kdy server DNS není dostupný. Soubor {\tt /etc/hosts} obvykle obsahuje jen základní mapování lokální adresy {\tt localhost} na IP adresu. Do souboru ale můžeme vložit i další doménová jména a odpovídající IP adresy, viz následující ukázka. Formát souboru lze najít v manuálových stránkách {\tt man hosts}. 
\begin{verbatim}
::1         localhost localhost.my.domain
127.0.0.1   localhost localhost.my.domain
10.0.0.3    myfriend.my.domain myfriend
\end{verbatim}
  \item {\tt /etc/resolv.conf} -- Jedná se o konfigurační soubor resolveru DNS. Resolver DNS vytváří dynamické mapování (nejen) IP adres a doménových jmen zasíláním dotazovaných jmen na server DNS. Konfigurační soubor obsahuje například název implicitní domény, které se doplňují za doménová jména, která nejsou plně kvalifikovaná (klíčové slovo {\tt search}), např. doménové jméno {\tt merlin} se při dotazování DNS automaticky rozšíří na {\tt merlin.fit.vutbr.cz}. Dále obsahuje konfigurační soubor {\tt /etc/resolv.conf} seznam serverů DNS (klíčové slovo {\tt nameserver}), na které lokální počítač posílají dotazy DNS. Popiš všech dostupných voleb konfigurace lze najít na manuálových stránkách {\tt man resolv.conf}. Konfigurace souboru je buď manuální, tj. administrátor ručně nakonfiguruje IP adresy lokálních serverů DNS, implicitní doménu apod., nebo se konfigurace doplní z odpovědi DHCP při počáteční konfiguraci systému. Pokud je nastavena konfigurace DHCP, tak se manuálně vložené data přepíšou informacemi z DHCP při nejbližší aktualizaci dat DHCP. Následující ukázka obsahuje příklad konfigurace lokálního klienta DNS na FIT:
\begin{verbatim}
# Generated by resolvconf
search fit.vutbr.cz
nameserver 147.229.9.43
nameserver 147.229.8.43
\end{verbatim}
  \item {\tt /etc/host.conf} -- U rezoluce doménových jmen je potřeba stanovit, zda má přednost statická konfigurace uložená v souboru {\tt /etc/hosts} nebo dynamická konfigurace přes {\tt resolver}. Pořadí rezoluce definuje konfigurační soubor {\tt /etc/host.conf}, který je automaticky generován z konfiguračního souboru {\tt nsswitch.conf}, viz {\tt man nsswitch.conf}. Soubor {\tt /etc/host.conf} udává pořadí zdrojů překladu DNS. Standardně se nejprve využívá statická konfigurace ze souboru {\tt /etc/hosts} a teprve když se dané doménové jméno nenajde, dotazuje se lokální resolver serveru DNS, viz následující ukázka konfigurace:
\begin{verbatim}
# Auto-generated from nsswitch.conf
hosts
dns
\end{verbatim}
\end{itemize}

\subsection{Konfigurace serveru DNS}
Konfiguraci serveru DNS tvoří jednak konfigurace vlastního serveru DNS (např. program {\tt named}, což je implemetace serveru Bind) a dále zónové soubory, které obsahují seznam záznamů DNS pro domény, které server spravuje. Konfigurace serveru DNS je uložena v konfiguračním souboru {\tt /etc/named.conf}, zónové soubory se ukládají do podadresář {\tt /var/named/}. 

\subsubsection{Konfigurační soubor {\tt named.conf}}
Jedná se o základní konfigurační soubor aplikace {\tt named}, viz {\tt man named.conf}. Tento soubor specifikuje, jak má lokální server DNS zpracovat přijatý dotaz DNS. DNS server se nejprve pokusí vyhledat dotazovanou doménu v konfiguračním souboru {\tt /etc/named.conf} v sekcích {\tt zone}. Potom začne prohledávat příslušný zónový soubor obsahující všechny platné záznamy v dané doméně. 

Pokud se jedná o lokálně spravovanou doménu, odpověď se hledá v lokálních zónových souborech umístěných na daném serveru v adresáři {\tt /var/named/}. Pokud server nenajde hledanou doménu v sekci {\tt zone} konfiguračního souboru {\tt /etc/named.conf}, znamená to, že daná doména není lokální. V takovém případě se pokusí server poslat dotaz (v případě rekurzivního vyhledávání) na kořenový server DNS ležící ve speciální zóně ".". Tato zóna má typ {\tt hint} a je také obvykle uvedena v konfiguračním souboru DNS serveru.

Zónový soubor kořenové domény mívá standardizované jméno {\tt named.root} (někdy také {\tt root.zone}) a obsahuje záznamy typu A a AAAA všech třinácti kořenových serverů DNS. Pokud zóna "." v konfiguračním souboru chybí a dotazovaná doména není lokální, vrátí server negativní odpověď.

Níže uvedený příklad obsahuje dvě lokální domény se jmény {\tt moje-domena.cz} a {\tt 10.10.10.in-addr .arpa}. U každé domény je odkaz na příslušný zónový soubor, kde jsou jednotlivé záznamy dané domény uloženy. Zároveň konfigurace server DNS obsahuje odkaz na kořenové servery (zóna typu {\tt hint}). 
{\footnotesize
\begin{verbatim}
...
zone "." IN {                        # zóna "." typu hint odkazující na kořenové servery DNS
  type hint;
  file "/etc/named/named.root";      # seznam kořenových serverů
};                                   # viz též https://www.internic.net/domain/named.root

zone "moje-domena.cz" IN {           # lokální doména se jménem "moje-domena.cz"
  type master;
  file "/var/named/moje-domena.cz";  # název zónového souboru obsahující záznamy dané domény
  notify no;
};

zone "10.10.10.in-addr.arpa" IN {    # lokální reverzní doména se jménem "10.10.10.in-addr.arpa"
  type master;
  file "/etc/named/10.10.10";        # název zónového souboru obsahující reverzní doménu
  notify no;
};

\end{verbatim}
}
%
%Zóna typu {\tt hint} odkazuje na seznam tzv. root serverů, tedy takových serverů, které znají informace
%o všech registrovaných top-level doménách a dokáží říci, které servery jsou za tyto top-level domény
%zodpovědné. Zóna typu {\tt master} značí, že server je autoritativním serverem pro danou doménu.
%Mezi povinné záznamy pro takovou doménu patří záznam typu SOA a alespoň jeden NS odkazující na server samotný. Zároveň by měl být obsažen i záznam typu A pro tento server.
%
%\subsubsection{Zóna typu hint}
%Asociuje se s~nejvyšší doménou v~rámci celé hierarchie. Pokud přijde na server požadavek, prohledá seznam spravovaných zón ostatních typů ({\em master}, {\em slave}, {\em forward}). Pokud záznam nenajde, vybere jeden z~kořenových serverů definovaný právě v~tomto typu zóny.
%
Soubor {\tt named.root}, který obsahuje záznamy A a AAAA kořenových serverů, bývá součastí instalace serveru DNS. Po instalaci dojde k aktualizaci tohoto seznamu podle veřejného seznamu InterNIC. Soubor {\tt named.root} lze také získat například programem dig:
\begin{verbatim}
 dig +norec NS . @a.root-servers.net > /var/named/db.root
\end{verbatim}

%% \subsubsection{Lokálně obsluhovane zóny}
%% Tyto zóny lze rozdělit do dvou skupin -- {\bf loopback adresy} a {\bf ,,prázdné zóny''}. RFC 6303
%% \cite{rfc6303} je z~kategorie {\em best practice} a definuje, které zóny by měl být schopen DNS server
%% obsloužit sám a tím redukovat dotazy přeposílané na další servery. Ve většině případů jsou to zóny pro reverzní záznamy privátních IP adres.
%% 
%Výchozí soubor aplikace {\bf named} ve FreeBSD respektuje právě RFC 6303 a navíc doplňuje ještě další
%zóny. Například nealokované rozsahy IPv6 adres. Pro práci v~laboratoři nahraďte původní soubor {\tt
%/etc/namedb/named.conf} souborem {\tt /root/isa1/named.conf}, který je zkrácen a obsahuje pouze nezbytné definice.

\subsubsection{Vytvoření vlastní domény}\label{sec:vlastni_zona}
Pro vytvoření vlastní domény je potřeba přidat do konfiguračního souboru {\tt /etc/named.conf} novou sekci {\tt zone} se jménem vlastní domény a typem {\tt master}. Dále je potřeba pro tuto doménu vytvořit zónový soubor, který bude obsahovat všechny potřebné záznamy DNS pro danou domény. Povinné záznamy jsou SOA a NS, volitelné pak A, AAAA, MX, CNAME a další. Pokud přidáváme  zónový soubor pro svou doménu, je vhodné vytvořit i reverzní zónový soubor pro zpětný překlad IP adres na doménová jména. Tento reverzní zónový soubor kromě povinných záznamů SOA a NS obsahuje dále jen záznamy typu PTR. Počet PTR záznamů by měl odpovídat počtu záznamů A a AAAA v přímém zónovém soub
oru, tj. aby pro každou IP adresu bylo možné najít příslušné doménové jméno. 

Příklad konfiguračního souboru DNS pro nově vytvářenou doménu {\tt isa.cz}:
\begin{verbatim}
zone "isa.cz" {
  type master;
  file "/var/named/isa.cz";
};
\end{verbatim}

Jednotlivé záznamy v doméně {\tt isa.cz} budou uloženy v souboru {\tt /var/named/isa.cz}. Název souboru může být libovolný (např. {\tt isa.cz}, {\tt cz.isa}, {\tt db.isa.cz}, {\tt isa.cz.zone}), záleží na zvyklostech daném administrátora DNS. Název zónového souboru nemá vliv na jméno vytvářené domény. Podstatné je uvést správné jméno domény za klíčové slovo {\tt zone}. Zde budeme pro jméno přímého zónového souboru používat jméno domény, tj. doména {\tt isa.cz} bude uložena v souboru {\tt isa.cz}. 

Formát zónového souboru obvykle obsahuje na začátku implicitní dobu životnosti dané zóny (klíčové slovo {\tt TTL}) a dále pak seznam jednotlivých záznamů pro danou doménu. Při vytváření záznamů je potřeba mít na paměti několik zásad:
\begin{itemize}
  \item Bílé znaky (prázdné řádky, odsazení) nejsou v konfiguraci důležité.
  \item Znak "@" má speciální význam - nahrazuje jméno domény, které je uvedeno v konfiguračním souboru za klíčovým slovem {\tt zone}. Tj. v našem případě se každý výskyt znaku "@" v zónovém souboru nahradí řetězce "isa.cz".
  \item Základní formát záznamu DNS je následující
\begin{verbatim}
      <domain-name> <TTL> <CLASS> <TYPE> <rdata>
např.   isa.cz.             IN      NS   ns.isa.cz.
\end{verbatim}
  \begin{itemize}
    \item \verb|<domain-name>| obsahuje doménové jméno, které se vyhledává. Toto jméno musí být tzv. plně kvalifikované (FQDN, Fully Qualified Domain Name), tedy musí obsahovat cestu v doménovém stromě až ke kořeni (včetně ukončení tečkou). Pokud doménové jméno není ukončení tečkou, automaticky se za ní přidá doména uvedená v názvu zóny. Pokud bychom například napsali pouze "isa.cz", pak DNS server doplní jména na "isa.cz.isa.cz.", což je asi něco jiného, než jsme požadovali. Naopak, pokud zapíše pouze název server "ns", pak server doplní název domény na "ns.isa.cz.", což je v pořádku.
    
    Pro záznamy typu {\tt SOA, NS, MX}, které popisují vlastnosti domény (tj. celého podstromu ve jmenném prostoru DNS), obsahuje toto políčko jméno dané domény, např. "isa.cz.". Pro záznamy záznamy typu {\tt A, AAA, CNAME}, které popisují vlastnosti konkrétního koncového uzlu jmenného prostoru DNS (například IP adresu, alias), obsahuje toto políčko doménové jméno (hostname) daného zařízení, např. "ns.isa.cz.", "pc1.isa.cz.", "www.isa.cz." apod. 
    
    \item \verb|<TTL>| udává dobu platnosti záznamu. Pokud tato hodnota není v popisu záznamu uvedena (jako v našem případě), vezme se implicitní hodnota uvedená na začátku zónového souboru za klíčovým slovem {\tt TTL}.
    \item \verb|<CLASS>| popisuje třídu záznamů DNS. Historicky existovalo více různých tříd záznamů DNS, v současné době se používá pouze jediná třída typu {\tt IN} (Internet). Tuto hodnotu je musí záznam DNS vždy obsahovat.
    \item \verb|<TYPE>| obsahuje typ záznamu DNS, například {\tt SOA, NS, A, MX, CNAME, PTR} apod. Typ záznamu také určuje, jaká hodnota bude na pravé straně záznamu.
    \item \verb|<rdata>| obsahuje data daného záznamu, jejichž formát je dán ve specifikace daného typu záznamu. Například záznam typu NS očekává na pravé straně doménové jméno serveru DNS pro danou doménu. Naopak záznam typu SOA obsahuje různé parametry dané zóny (např. jméno autoritativního serveru DNS, e-mailovou adresu správce, hodnoty TTL pro aktualizaci záznamů apod.
  \end{itemize}
  \item Znak ";" se používá jako oddělovač komentářů
\end{itemize}

\subsubsection{Příklady záznamů DNS v přímém zónovém souboru}
Jak již bylo řečeno, každý zónový soubor obsahuje právě jeden záznam SOA s údaji o dané zóně a minimálně jeden záznam typu NS s adresou autoritativního serveru DNS pro danou doménu. Kromě těchto záznamů obvykle obsahuje zónový soubor záznamy typu A (překlad doménových jmen na IPv4 adresy), záznamy typu AAAA (překlad doménových jmen na adresa IPv6), záznamy typu MX (odkazujících na poštovní server pro danou doménu), záznamy typu CNAME (aliasy koncových zařízení) a další. 

\begin{itemize}
  \item SOA -- základní informace o zóně
\begin{verbatim}
@ IN  SOA ns.isa.cz. admin.isa.cz. ( ; primární server DNS, e-mail na správce
          2023092501   ; seriové číslo - datum editace, pořadí 01
          1w  ; refresh: interval obnovy pro sekundární servery (1 týden)
          1d  ; retry: po jaké době se pokusí sekundární server kontaktovat
              ; primární server, pokud se předchozi spojení nezdařilo (1 den)
          1w  ; expire: po jaké době je záznam neplatný, pokud se nepodařila
              ; aktualizace
          1h  ; negative cache TTL: TTL pro chybové odpovědi (např. NXDOMAIN)
  )
\end{verbatim}
  \clearpage
  \item NS -- autoritativní servery DNS 
\begin{verbatim}
  @       IN NS ns1.isa.cz.  ; zápis využívající zástupný symbol "@"
  isa.cz. IN NS ns1.isa.cz.  ; ekvivalentní zápis s plným doménovým jménem
  @     IN NS ns1            ; zkrácený ekvivalentní zápis
\end{verbatim}        
  \item MX -- poštovní server pro danou doménu (s prioritou)
\begin{verbatim}
  isa.cz. IN MX 10 mail.isa.cz. ; primární poštovní server 
  isa.cz. IN MX 20 mail2.isa.cz.; sekundární poštovní server
  isa.cz. IN MX 20 mail2        ; zkrácený ekvivalentní zápis
  @       IN MX 20 mail2        ; zkrácený ekvivalentní zápis
\end{verbatim}        
  \item A -- překlad IPv4 adres na doménová jména
\begin{verbatim}
  ns.isa.cz. IN A 192.168.1.1  ; plně kvalifikované doménové jméno
  ns         IN A 192.168.1.1  ; zkrácený ekvivalentní zápis
  pc1        IN A 192.168.1.2  ; IP adresa pro pc1.isa.cz.
  pc2        IN A 192.168.1.3  ; IP adresa pro pc2.isa.cz.
\end{verbatim}        
  \item AAAA -- překlad adres IPv6 pro doménová jména
\begin{verbatim}
  pc3        IN AAAA fd12:1234::1  ; adresa IPv6 pro pc3.isa.cz.
\end{verbatim}        
  \item CNAME -- alias počítače
\begin{verbatim}
  www.isa.cz.  IN CNAME pc1.isa.cz. ; alias www pro PC1
  www          IN CNAME pc1         ; zkrácený ekvivalentní zápis
  mail         IN CNAME pc1         ; alias pro server mail
  mail2        IN CNAME pc2         ; alias pro server mail2   
\end{verbatim}        
\end{itemize}

\subsubsection{Reverzní zónový soubor}
Reverzní zónový soubor obsahuje záznamy typu {\tt PTR}, které překládají IP adresu na doménové jméno. Pravá strana záznamu leží ve speciální podstromu {\tt in-addr.arpa.} pro adresy IPv4 nebo v podstromu {\tt ip6.arpa} pro IPv6. Záznamy typu {\tt PTR} se vytvářejí pro každý přímý překlad, tj. pro každý záznam {\tt A} nebo {\tt AAAA} by měl existovat odpovídající reverzní záznam. V našem příkladu máme čtyři přímé záznamy pro doménová jména {\tt ns.isa.cz, pc1.isa.cz, pc2.isa.cz} a {\tt pc3.isa.cz}, pro které budeme vytvářet reverzní záznamy. 

Reverzní záznamy jsou uloženy v jiném zónovém souboru než přímé záznamy. Pro název přímého zónového souboru jsme použili jméno domény (např. "isa.cz"), pro název reverzního zónového souboru se obvykle používá prefix podsítě, do které patří počítače v dané doméně. Tento prefix se zapisuje v opačné pořadí bytů. Tedy pro podsíť IPv4 {\tt 192.168.1.0} bude název reverzní zóny "1.168.192.in-addr.arpa". Tento název se objeví i v konfiguraci serveru DNS v souboru {\tt /etc/named.conf}, kde bude potřeba vytvořit novou zónu, viz následující příklad:
\begin{verbatim}
zone "1.168.192.in-addr.arpa" {
  type master;
  file "/var/named/192.168.1";
};
\end{verbatim}        
Reverzní zónový soubor jsme v tomto případě pojmenovali {\tt 192.168.1}, nicméně název opět závisí na zvyklostech administrátora DNS. Podobně jako přímý zónový soubor musí reverzní zónový soubor obsahovat dobu platnosti zóny {\tt TTL}, dále povinné záznamy SOA a NS. Dále obsahuje záznamy PTR.

Příklad reverzního zónového souboru pro doménu {\tt isa.cz} je v následující ukázce:
\begin{verbatim}
$TTL 3h

@        IN SOA ns.isa.cz. admin.isa.cz. (            ; SOA záznam
               2023092501 ; serial number
               1w         ; refresh
               1d         ; retry 
               1w         ; expire
               1h         ; negative cache TTL
         )
@                         IN NS ns.isa.cz.  ; NS záznam (zkrácený)
1.168.192.in-addr.arpa.   IN NS ns.isa.cz.  ; plný zápis

1.1.168.192.in-addr.arpa. IN PTR ns.isa.cz.  ; plný zápis pro ns
1                         IN PTR ns.isa.cz.  ; zkrácený zápis pro ns
2                         IN PTR pc1.isa.cz. ; zrkácený zápis pro pc1
3                         IN PTR pc2.isa.cz. ; zkrácený zápis pro pc2
                    
\end{verbatim}
Při zápisu je potřeba si uvědomit, že na pravé straně musí být vždy plně kvalifikované doménové jméno. Pokud bychom na pravé straně uvedli například jenom "ns", tak po vyhodnocení se přidá "ns.1.168.192.in-addr.arpa.", což asi nechceme. Podobně když zapomeneme tečku, tj. napíšeme pouze "ns.isa.cz", vyhodnocení opět dopadne jinak, než si představujeme -- server vrátí hodnotu "ns.isa.cz.1 .168.192.in-addr.arpa.". 

Pro reverzní překlad IPv6 bude muset existovat nový zónový soubor, protože reverzní adresy IPv6 mají jiný podstrom v DNS než adresy IPv4. Do konfiguračního souboru {\tt /etc/named.conf} přidáme novou zónu, např. "2.1.d.f.ip6.arpa.", viz následující ukázka:
\begin{verbatim}
zone "2.1.d.f.ip6.arpa" {
  type master;
  file "/var/named/fd12";
};
\end{verbatim} 
Tento reverzní zónový soubor pro zpětný překlad adres IPv6 na doménová jména bude opět obsahovat záznam SOA a NS, dále pak záznamy PTR v doméně  "2.1.d.f.ip6.arpa", viz následující ukázka: 
\begin{verbatim}
$TTL 3h

@        IN SOA ns.isa.cz. admin.isa.cz. (            ; SOA záznam
               2023092501 ; serial number
               1w         ; refresh
               1d         ; retry 
               1w         ; expire
               1h         ; negative cache TTL
         )

@                   IN NS ns.isa.cz.  ; NS záznam (zkrácený)
2.1.d.f.ip6.arpa.   IN NS ns.isa.cz.  ; plný zápis

1.0.0.0.0.0.0.0.0.0.0.0.0.0.0.0.0.0.0.0.0.0.0.0.4.3.2.1.2.1.d.f.ip6.arpa. IN PTR pc3.isa.cz.
                                      ; plný zápis pro adresu fd12:1234::1
1.0.0.0.0.0.0.0.0.0.0.0.0.0.0.0.0.0.0.0.0.0.0.0.4.3.2.1 IN PTR pc3.isa.cz.
                                      ; zkrácený zápis pro adresu fd12:1234::1
\end{verbatim}
Nevýhodou reverzního zápis adresy IPv6 je, že záznamy DNS neznají zkracování nulových bytů. Je nutné uvést celou cestu v reverzním stromu DNS pro IPv6, kde každá úroveň má 8 bitů, tj. jeden byte. Adresa IPv6 obsahuje 32 bytů, které je nutné v reverzním překladu uvést. Protože tento výpis je značně složitý a může dojít snad k chybě, používají se na vytvoření těchto adres specializované generátory\footnote{Viz například \url{https://www.whatsmydns.net/reverse-dns-generator}.}. 

Po každé změně zónového souboru by se vždy mělo upravit i sériové číslo zóny a to tak, aby bylo vždy větší než předchozí hodnota. Také je nutné aplikaci načíst změněnou konfiguraci.
\begin{verbatim}
systemctl restart named.service
\end{verbatim}

%% \subsection{DNS Over TLS}\label{DoT}
%% DNS Over TLS (DOT) přináší šifrování a autentizaci do DNS komunikace. Tento protokol je definován v RFC 7858\cite{RFC7858}. Protokol má přidělen standardizovaný port 853 jako defaultní. TLS protokol zde zaobaluje DNS požadavky a odpovědi jako přenášený payload viz obrázek \ref{fig:dot_encapsulation}.
%% 
%% \begin{figure}[h]
%%     \centering
%%     \includegraphics[scale=1]{fig/dot.pdf}
%%     \caption{DNS Over TLS zapouzdření.}
%%     \label{fig:dot_encapsulation}
%% \end{figure}
%% 
%% Komunikace probíhá pomocí klasického klient-server schématu a transportním protokolem je TCP. Nejprve proběhne TCP handshake, následuje TLS handshake tak jak je definován v RFC 5246\cite{RFC5246}. Jakmile je spojení ustaveno dojde k přenosu dat.
%% 
%% Pro minimalizaci latence klienti nemusí čekat na odpovědi serveru, ale můžou použít pipelining pro odesílání více dotazů v rámci jednoho TLS spojení. Přenášená data uvnitř TLS jsou stejného formátu jako v případě standardního DNS přenášeného přes TCP transportní protokol.\cite{RFC7858}
%% 
%% Resolvíng pomocí DNS over TLS lze na Linuxových systémech nastavit například v podobě spuštění Unbound DNS Caching serveru. Kde je nutné upravit konfigurační soubor \texttt{/etc/unbound/unbound.conf} tak, aby obsahoval následující:
%% 
%% \begin{verbatim}
%% forward-zone:
%%     name: "."
%%     forward-ssl-upstream: yes
%%     # Cloudflare DNS
%%     forward-addr: 1.1.1.1@853
%%     forward-addr: 1.0.0.1@853
%% \end{verbatim}
%% 
%% Následně službu zapnout a nastavit systémový resolver na adresu \texttt{127.0.0.1}.
%% 
%% \subsection{DNS Over HTTPS}\label{DoH}
%% DNS Over HTTPS (DOH) je nový standard pro bezpečný a šifrovanz přenos DNS provozu tunelovaný v protokolu HTTPS. Tento protokol nemá vlastní definovaný port a je přenášen jako ostatní HTTPS provoz pod portem 443. DOH je definován standardem vydaným roku 2018 skrývajícím se pod dokumentem označeným jako RFC 8484 \cite{RFC8484}.
%% 
%% \begin{figure}[h]
%%     \centering
%%     \includegraphics[scale=0.75]{fig/doh-encapsulation.pdf}
%%     \caption{DNS over HTTPS zapouzdření.}
%%     \label{fig:doh_encapsulation}
%% \end{figure}
%% 
%% DNS zprávy jsou zde přenášeny standardními metodami HTTP GET a POST. V případě GET je DNS request zakódován ve formátu URL64Base jako parametr URI. Výsledná odpověď je obdržena v těle HTTP odpovědi ve standardním DNS wire formátu. V případě metody POST je již sám dotaz zakódován v DNS wire formátu a přenášen v těle POST dotazu.
%% 
